
\section{\BookRecordIndex{T0145}格罗莫夫--威滕伪循环}\label{sec:6-6}
\zhindex{geluomofuweitengweixunhuan}{格罗莫夫--威滕伪循环}

现在说明本章的主要结果：在适当条件下，\BookRecordIndex{T0283}求值映射
\[
  \operatorname{ev}:\mathcal M^*_{0,k}(A;J)\longrightarrow M^k
\]
是\BookRecordIndex{T0337}伪循环。当 $A=mB$、$m>1$，而 $c_1(B)=0$ 时，这里给出的论证并不适用。处理这种
情形的最佳方法，是用曲线 $u:S^2\to M$ 的图像代替曲线本身；下一节将讨论这一方法。

回忆第~\ref{sec:6-2}节中集合
\[
  \mathcal J_{\mathrm{reg}}(M,\omega)
  :=\bigcap_T\ \bigcap_{\{A_\alpha\}}
    \mathcal J_{\mathrm{reg}}(T,\{A_\alpha\})
\]
的定义。它由 $M$ 上所有被 $\omega$ \BookRecordIndex{T0372}驯服并满足下述条件的\BookRecordIndex{T0088}殆复结构 $J$ 组成：对每个
单纯 $J$-全纯球面 $u:\mathbb{CP}^1\to M$，$D_u$ 都满射；而且对每棵\BookRecordIndex{T0084}带标号树 $T$
及每族球面同调类 $\{A_\alpha\}$，\BookRecordIndex{T0283}求值映射
\[
  \operatorname{ev}^E:
  \mathcal M^*(\{A_\alpha\};J)\times Z(T)
  \longrightarrow M^E
\]
都横截于 $\Delta^E$。由定理~\ref{thm:6-2-6}，
$\mathcal J_{\mathrm{reg}}(M,\omega)$ 在
$\mathcal J_\tau(M,\omega)$ 中是剩余的。

\begin{theorem}\label{thm:6-6-1}
设 $(M,\omega)$ 是闭\BookRecordIndex{T0058}半正辛 $2n$-流形，且
$J\in\mathcal J_{\mathrm{reg}}(M,\omega)$。设 $A\in H_2(M;\mathbb Z)$ 满足：对每个
正整数 $m$ 和每个球面同调类 $B\in H_2(M;\mathbb Z)$，
\begin{equation}\label{eq:6-6-1}
  A=mB,\qquad c_1(B)=0
  \qquad\Longrightarrow\qquad m=1.
\end{equation}
则\BookRecordIndex{T0283}求值映射
\[
  \operatorname{ev}:\mathcal M^*_{0,k}(A;J)\longrightarrow M^k
\]
是维数
\[
  \mu(A,k)=2n+2c_1(A)+2k-6
\]
的\BookRecordIndex{T0337}伪循环，而且它的配边类与 $J$ 的选择无关。
\end{theorem}

\begin{proof}
\BookRecordIndex{T0139}格罗莫夫紧性定理（定理~\ref{thm:5-3-1}）蕴含：\BookRecordIndex{T0283}求值映射
$\operatorname{ev}$ 的极限集
\[
  \Omega_{\operatorname{ev}}
  :=
  \bigcap_{\substack{K\subset\mathcal M^*_{0,k}(A;J)\\K\text{ 紧}}}
  \overline{\operatorname{ev}
  \bigl(\mathcal M^*_{0,k}(A;J)\setminus K\bigr)}
\]
由下述各像之并覆盖：其一，是至少带一条边的各棵 $k$-\BookRecordIndex{T0084}带标号树 $T$ 所对应的\BookRecordIndex{T0072}层
$\mathcal M_{0,T}(A;J)$ 在\BookRecordIndex{T0283}求值映射
\sourcefootnote{原书此处写成“至少两条边”；但一条边的\BookRecordIndex{T0340}稳定映射同样是主\BookRecordIndex{T0072}层的边界，
而下文的维数论证明确对所有 $e(T)>0$ 的树进行。故唯一补为“至少一条边”。}
\[
  \operatorname{ev}_T:\mathcal M_{0,T}(A;J)\longrightarrow M^k,
  \qquad
  \operatorname{ev}_T(\mathbf u,\mathbf z)
  :=\bigl(u_{\alpha_1}(z_1),\ldots,u_{\alpha_k}(z_k)\bigr)
\]
下的像；其二，是\BookRecordIndex{T0113}多重覆盖曲线空间
$\mathcal M_{0,k}(A;J)\setminus\mathcal M^*_{0,k}(A;J)$ 在
$\operatorname{ev}$ 下的像。

接着注意，由命题~\ref{prop:6-1-2}，只要扩大所容许的表示同调类的范围，就可以用\BookRecordIndex{T0184}简单
稳定映射的像覆盖 $\Omega_{\operatorname{ev}}$。具体地，
\begin{equation}\label{eq:6-6-2}
  \Omega_{\operatorname{ev}}
  \subset
  \bigcup_T\ \bigcup_{\{B_\alpha\}}
  \operatorname{ev}_T
  \bigl(\mathcal M^*_{0,T}(\{B_\alpha\};J)\bigr),
\end{equation}
其中并集遍历所有 $k$-\BookRecordIndex{T0084}带标号树 $T$ 以及所有满足下列条件的球面同调类族
$\{B_\alpha\}_{\alpha\in T}$：
\begin{equation}\label{eq:6-6-3}
  A=\sum_{\alpha\in T}m_\alpha B_\alpha
\end{equation}
以及
\begin{equation}\label{eq:6-6-4}
  e(T)=0\qquad\Longrightarrow\qquad B\ne A,
\end{equation}
这里 $\{m_\alpha\}_{\alpha\in T}$ 是某族正整数。

这一公式对每个紧\BookRecordIndex{T0358}辛流形 $(M,\omega)$ 和每个\BookRecordIndex{T0088}殆复结构
$J\in\mathcal J(M,\omega)$ 都成立。由于
$J\in\mathcal J_{\mathrm{reg}}(M,\omega)$，定理~\ref{thm:6-2-6} 表明\BookRecordIndex{T0246}模空间
$\mathcal M^*_{0,T}(\{B_\alpha\};J)$ 是光滑流形，其维数为
\[
  \dim\mathcal M^*_{0,T}(\{B_\alpha\};J)
  =\mu(B,k)-2e(T),
  \qquad
  B:=\sum_{\alpha\in T}B_\alpha.
\]
由引理~\ref{lem:6-4-4}，\BookRecordIndex{T0058}半正性假设保证
$\mathcal M^*_{0,T}(\{B_\alpha\};J)$ 为空，除非对每个 $\alpha\in T$ 都有
\[
  c_1(B_\alpha)\ge0.
\]
由式~\eqref{eq:6-6-3}，对每个非空的\BookRecordIndex{T0072}层，这蕴含
$\mu(B,k)\le\mu(A,k)$。所以，对每棵满足 $e(T)>0$ 的树，都有
\[
  \mu(B,k)-2e(T)\le\mu(A,k)-2.
\]
若 $e(T)=0$，则由式~\eqref{eq:6-6-3} 和~\eqref{eq:6-6-4}，有
$A=mB$，其中 $m\ge2$ 是整数。由式~\eqref{eq:6-6-1}，这蕴含
$c_1(B)\ne0$；再由\BookRecordIndex{T0058}半正性，$c_1(B)>0$，所以
\[
  \mu(B,k)\le\mu(A,k)-2.
\]
这就证明了
$\operatorname{ev}:\mathcal M^*_{0,k}(A;J)\to M^k$ 是\BookRecordIndex{T0337}伪循环。

现在取 $J_0,J_1\in\mathcal J_{\mathrm{reg}}(M,\omega)$，并考虑相应的\BookRecordIndex{T0283}求值映射
\[
  \operatorname{ev}_0:\mathcal M^*_{0,k}(A;J_0)\longrightarrow M^k,
  \qquad
  \operatorname{ev}_1:\mathcal M^*_{0,k}(A;J_1)\longrightarrow M^k.
\]
我们证明，这两个映射在定义~\ref{def:6-5-1} 的意义下作为\BookRecordIndex{T0337}伪循环彼此配边。回忆，
$\mathcal J_{\mathrm{reg}}(M,\omega;J_0,J_1)$ 表示所有连接 $J_0,J_1$ 的光滑同伦
$\{J_\lambda\}_{0\le\lambda\le1}$ 所成的集合，其中同伦对每棵树 $T$ 和每族同调类
$\{A_\alpha\}_{\alpha\in T}$ 都在定义~\ref{def:6-2-10} 的意义下正则。由
定理~\ref{thm:6-2-11}，集合
$\mathcal J_{\mathrm{reg}}(M,\omega;J_0,J_1)$ 非空。固定\BookRecordIndex{T0407}正则同伦
\[
  \{J_\lambda\}\in
  \mathcal J_{\mathrm{reg}}(M,\omega;J_0,J_1),
\]
并考虑\BookRecordIndex{T0283}求值映射
\[
  \operatorname{ev}:
  \mathcal W^*_{0,k}(A;\{J_\lambda\})
  \longrightarrow M^k.
\]
由定理~\ref{thm:6-2-11}，空间
$\mathcal W^*_{0,k}(A;\{J_\lambda\})$ 是定向光滑带边流形，并且
\[
  \partial\mathcal W^*_{0,k}(A;\{J_\lambda\})
  =
  \mathcal M^*_{0,k}(A;J_1)
  \sqcup\bigl(-\mathcal M^*_{0,k}(A;J_0)\bigr).
\]
它的维数为
\[
  \dim\mathcal W^*_{0,k}(A;\{J_\lambda\})
  =\mu(A,k)+1,
\]
而\BookRecordIndex{T0283}求值映射在 $\lambda=0$ 与 $\lambda=1$ 两个边界部分上的限制，显然分别是
$\operatorname{ev}_0$ 与 $\operatorname{ev}_1$。

再次由定理~\ref{thm:5-3-1} 和命题~\ref{prop:6-1-2}，
\[
  \bigcap_{\substack{
    K\subset\mathcal W^*_{0,k}(A;\{J_\lambda\})\\K\text{ 紧}}}
  \overline{\operatorname{ev}
  \bigl(\mathcal W^*_{0,k}(A;\{J_\lambda\})\setminus K\bigr)}
  \subset
  \bigcup_T\ \bigcup_{\{B_\alpha\}}
  \operatorname{ev}_T
  \bigl(\mathcal W^*_{0,T}(\{B_\alpha\};\{J_\lambda\})\bigr),
\]
其中并集遍历所有 $k$-\BookRecordIndex{T0084}带标号树 $T$ 和所有满足式~\eqref{eq:6-6-3}
与~\eqref{eq:6-6-4} 的球面同调类族 $\{B_\alpha\}_{\alpha\in T}$。由
定理~\ref{thm:6-2-11}，这个并集中的\BookRecordIndex{T0246}模空间维数为
\[
  \dim\mathcal W^*_{0,T}(\{B_\alpha\};\{J_\lambda\})
  =\mu(B,k)+1-2e(T)
  \le\mu(A,k)-1.
\]
最后一个不等式来自证明第一部分中的同一论证。注意，这里必须把
引理~\ref{lem:6-4-4} 用于未必正则的 $J_\lambda$；不过，它们满足
$\dim\operatorname{coker}D_u\le1$。因此，
\[
  \operatorname{ev}:
  \mathcal W^*_{0,k}(A;\{J_\lambda\})\longrightarrow M^k
\]
正是从 $\operatorname{ev}_0$ 到 $\operatorname{ev}_1$ 的所需配边。
\end{proof}

\begin{remark}\label{rem:6-6-2}
在非\BookRecordIndex{T0058}半正情形，只要仅考虑开集
\[
  \mathcal J_+(M,\omega;\kappa)
  \subset\mathcal J_\tau(M,\omega),
  \qquad \kappa=\omega(A),
\]
中的\BookRecordIndex{T0088}殆复结构，定理~\ref{thm:6-6-1} 的断言仍成立（见
引理~\ref{lem:6-4-7}）。在这种情形，对每个
\[
  J\in
  \mathcal J_{\mathrm{reg}}(M,\omega)
  \cap\mathcal J_+(M,\omega;\kappa),
\]
\BookRecordIndex{T0283}求值映射都是\BookRecordIndex{T0337}伪循环。不过，一般而言
$\mathcal J_+(M,\omega;\kappa)$ 可能为空；即使非空，也可能不连通。关于与 $J$ 选择
无关的断言，只对上述交集中属于 $\mathcal J_+(M,\omega;\kappa)$ 同一道路连通分支的
那些 $J$ 成立。
\end{remark}

\begin{exercise}\label{ex:6-6-3}
设 $(M,\omega)$ 是紧\BookRecordIndex{T0058}半正\BookRecordIndex{T0358}辛流形，
$J\in\mathcal J_{\mathrm{reg}}(M,\omega)$。设 $T$ 是 $k$-\BookRecordIndex{T0084}带标号树，
$\{A_\alpha\}_{\alpha\in T}$ 是满足稳定性条件~\eqref{eq:6-1-1} 的球面同调类族。
假设每个类 $A_\alpha$ 都满足式~\eqref{eq:6-6-1}。证明\BookRecordIndex{T0283}求值映射
\[
  \operatorname{ev}_T:
  \mathcal M^*_{0,T}(\{A_\alpha\};J)
  \longrightarrow M^k
\]
是维数为 $\mu(A,k)-2e(T)$ 的\BookRecordIndex{T0337}伪循环。
\end{exercise}

\paragraph{\BookRecordIndex{T0150}固定标记点。}
$S^2$ 上互异点的有限族
$\mathbf w:=(w_1,\ldots,w_k)$ 决定\BookRecordIndex{T0283}求值映射
\[
  \operatorname{ev}_{\mathbf w}:
  \mathcal M^*(A;J)\longrightarrow M^k.
\]
人们可能会预期，对一般的\BookRecordIndex{T0088}殆复结构，这是一个维数为
$2n+2c_1(A)$ 的\BookRecordIndex{T0337}伪循环。然而，除非对 $(M,\omega)$ 和\BookRecordIndex{T0066}标记点数施加相当强的限制，
事实似乎并非如此。如下例所示，问题在于：在含\BookRecordIndex{T0113}多重覆盖曲线的\BookRecordIndex{T0072}层上，会失去对\BookRecordIndex{T0066}标记点位置
的控制。

\begin{example}\label{ex:6-6-4}
设 $(M,\omega)$ 是一个含有类 $A$ 的六维流形，且 $c_1(A)=2$。例如，可以取
\[
  M=S^2\times\mathbb{CP}^2,\qquad A=[S^2\times\mathrm{pt}].
\]
考虑\BookRecordIndex{T0283}求值映射
\[
  \operatorname{ev}_{\mathbf w}:
  \mathcal M^*(2A;J)\longrightarrow M^5.
\]
完全可能发生如下情形：对某个正则 $J$，一列类 $2A$ 中的单纯 \BookRecordIndex{T0003}$J$-全纯曲线收敛到
一条双重覆盖的 $A$-曲线；见下文。因此，$\operatorname{ev}_{\mathbf w}$ 的极限集含有
形如 $u\circ\phi$ 的双重覆盖 $A$-曲线，其中 $u:S^2\to M$ 单纯，而
$\phi:S^2\to S^2$ 的度数为 $2$。$S^2$ 上二次全纯自映射的空间有复维数 $5$；用齐次
坐标可将它们写成
\[
  [z:w]\longmapsto
  [az^2+bzw+cw^2:dz^2+ezw+fw^2].
\]
容易验证，存在这样的映射把任意一个由五个互异点组成的集合 $\mathbf w$ 映到任意另一
个由五个互异点组成的集合。因此，若试图沿用定理~\ref{thm:6-6-1} 的证明，用
$u\circ\phi$ 的底\BookRecordIndex{T0072}层单纯映射 $u$ 代替它，为覆盖这个边界\BookRecordIndex{T0072}层，就必须允许求值点
$\mathbf w$ 在 $S^2$ 上自由变化。所以，这个\BookRecordIndex{T0072}层由定义域为
$\mathcal M_{0,5}(A;J)$ 的\BookRecordIndex{T0283}求值映射覆盖，而该空间与
$\mathcal M^*(2A;J)$ 维数相同。

为说明这个问题并不只是证明方法造成的假象，考虑如下例子。设 $L_{-1}$ 是
$\mathbb{CP}^2$ 上的线丛，其陈类在直线 $\ell$ 上取值 $-1$，并令
\[
  M:=\mathbb P(L_{-1}\oplus\mathbb C).
\]
那么 $M$ 含有一个\BookRecordIndex{T0117}法丛为 $L_{-1}$ 的 $\mathbb{CP}^2$ 副本 $Z$。若 $A$ 是 $Z$ 中
直线 $\ell$ 的类，则 $c_1(A)=2$。此外，类 $2A$ 中的邻近曲线是
$Z=\mathbb{CP}^2$ 中的圆锥曲线；可以验证，对 $\ell$ 上任意五个点
$x_1,\ldots,x_5$ 和任意 $\mathbf w$，都存在一列嵌入圆锥曲线
$u^\nu:S^2\to Z$，使得
\[
  u^\nu(w_i)\longrightarrow x_i,\qquad i=1,\ldots,5,
\]
当 $\nu\to\infty$。为看出这一点，选取收敛到 $x_i$ 的点 $y_i^\nu$，使它们对每个
$\nu$ 都落在一对接近 $\ell$ 的直线上，再把这对直线稍作扰动，得到非退化圆锥曲线。
\end{example}

注意，在上例中，尽管 $A$-曲线的\BookRecordIndex{T0246}模空间维数小于 $2A$-曲线的\BookRecordIndex{T0246}模空间维数（因为
$c_1(A)>0$），这种维数降低仍不足以补偿对\BookRecordIndex{T0066}标记点失去的控制。不过，若假设\BookRecordIndex{T0066}标记点数
$k$ 相对于\BookRecordIndex{T0418}最小陈数 $N$ 足够小，就能证明
$\operatorname{ev}_{\mathbf w}$ 是\BookRecordIndex{T0337}伪循环。下面的习题~\ref{ex:6-6-5} 给出一个示例
结果。这种解决方案显然不能令人满意。像本书第一版那样，通过更细致的讨论可以得到更强
的结果；但好得多的思路，是扩大所容许的 $J$ 的扰动范围，从而完全避开麻烦的\BookRecordIndex{T0113}多重覆盖
曲线。下一节采用的正是这种方法。

下一个例子对流形 $M$ 施加了很强的限制；作为交换，它适用于更一般的\BookRecordIndex{T0283}求值映射，其中只有
一部分\BookRecordIndex{T0066}标记点在定义域中固定。

\begin{exercise}\label{ex:6-6-5}
设 $(M,\omega)$ 是\BookRecordIndex{T0418}最小陈数为 $N$ 的单调\BookRecordIndex{T0358}辛流形，且
$J\in\mathcal J_{\mathrm{reg}}(M,\omega)$。固定同调类
$A\in H_2(M;\mathbb Z)$、整数 $k\ge3$、指标集
$I\subset\{1,\ldots,k\}$，并假设
\[
  3\le\#I\le N+2.
\]
再固定 $S^2$ 上互异点的有限族
\[
  \mathbf w:=\{w_i\}_{i\in I}.
\]
考虑所有元组 $(u,z_1,\ldots,z_k)$ 组成的\BookRecordIndex{T0246}模空间
$\mathcal M^*_{0,k}(A;\mathbf w,J)$，其中 $u:S^2\to M$ 是类 $A$ 中的单纯
$J$-全纯球面，$z_1,\ldots,z_k$ 是 $S^2$ 上两两互异的点，并且当 $i\in I$ 时
$z_i=w_i$。证明\BookRecordIndex{T0283}求值映射
\[
  \operatorname{ev}_{\mathbf w}:
  \mathcal M^*_{0,k}(A;\mathbf w,J)\longrightarrow M^k,
\]
定义为
\[
  \operatorname{ev}_{\mathbf w}(u,z_1,\ldots,z_k)
  :=\bigl(u(z_1),\ldots,u(z_k)\bigr),
\]
是维数为
\[
  2n+2c_1(A)+2(k-\#I)
\]
的\BookRecordIndex{T0337}伪循环。

\emph{提示：}设
$(u^\nu,z_1^\nu,\ldots,z_k^\nu)$ 是
$\mathcal M^*_{0,k}(A;\mathbf w,J)$ 中的一列，并且没有以
$\mathcal M^*_{0,k}(A;\mathbf w,J)$ 中元素为极限的收敛子列。由
定理~\ref{thm:5-3-1}，它仍有一个子列\BookRecordIndex{T0141}格罗莫夫收敛到某个\BookRecordIndex{T0340}稳定映射
$(\mathbf u,\mathbf z)$，其模型是一棵 $k$-\BookRecordIndex{T0084}带标号树 $(T,E,\Lambda)$。证明存在唯一的
\BookRecordIndex{T0321}特殊顶点 $0\in T$，其特征是诸点 $z_{0i}$（$i\in I$）两两互异。还要证明，与这个\BookRecordIndex{T0321}特殊
顶点相应的 $J$-全纯球面 $u_0:S^2\to M$ 可参数化为：对每个 $i\in I$，
$z_{0i}=w_i$。

接着考虑极限\BookRecordIndex{T0340}稳定映射单纯，而且边数满足
\[
  1\le e(T)\le\#I-3
\]
的情形。证明在这一假设下，\BookRecordIndex{T0321}特殊顶点承载指标集 $I$ 中至少
\[
  \#I-e(T)\ge3
\]
个\BookRecordIndex{T0066}标记点。证明：删去这些点中的 $\#I-e(T)-2$ 个，并允许其余点（至少两个）自由变化
后，所得映射仍然稳定，而且属于一个维数为
\[
  2n+2c_1(A)+2(k-\#I)-2
\]
的\BookRecordIndex{T0246}模空间。在所有其他情形，证明既约的\BookRecordIndex{T0184}简单稳定映射（带任意\BookRecordIndex{T0066}标记点）属于维数严格小于
$2n+2c_1(A)+2(k-\#I)$ 的\BookRecordIndex{T0246}模空间。
\end{exercise}

\begin{exercise}\label{ex:6-6-6}
设 $M=\mathbb{CP}^2$，$L=[\mathbb{CP}^1]$，而 $J_0$ 是 $M$ 上的标准复结构。
\begin{enumerate}[label=(\roman*),leftmargin=3em]
\item 证明非单纯元素
\[
  \mathcal M_{0,5}(2L;J_0)\setminus
  \mathcal M^*_{0,5}(2L;J_0)
\]
构成流形，并计算其维数。把它与底\BookRecordIndex{T0072}层\BookRecordIndex{T0183}简单层
$\mathcal M_{0,5}(L;J_0)$ 的维数比较。
\item 证明对每个 $k$，\BookRecordIndex{T0283}求值映射
\[
  \operatorname{ev}:\mathcal M^*_{0,k}(2L;J_0)\longrightarrow M^k
\]
都是\BookRecordIndex{T0337}伪循环。
\item 考虑习题~\ref{ex:6-6-5} 中 $k=|I|=4$ 的情形，并取
\[
  \mathbf w:=(0,1,\infty,-1).
\]
证明\BookRecordIndex{T0283}求值映射
\[
  \operatorname{ev}_{\mathbf w}:
  \mathcal M^*_{0,4}(2L;\mathbf w,J_0)
  \longrightarrow M^4
\]
是\BookRecordIndex{T0337}伪循环。注意，
$\operatorname{ev}_{\mathbf w}$ 是相同维数流形之间的映射。
\end{enumerate}
\end{exercise}
